\documentclass[11pt,a4paper]{article}
\usepackage[margin=25mm,headheight=14pt]{geometry}
\usepackage{amsmath,amssymb,amsthm,array,url,xcolor}
\usepackage[colorlinks=true,allcolors=inkblue,breaklinks=true]{hyperref}
\definecolor{inkblue}{RGB}{28,53,79}
\newtheorem{theorem}{Theorem}[section]
\newtheorem{lemma}[theorem]{Lemma}
\newtheorem{proposition}[theorem]{Proposition}
\newtheorem{corollary}[theorem]{Corollary}
\theoremstyle{definition}\newtheorem{definition}[theorem]{Definition}
\theoremstyle{remark}\newtheorem{remark}[theorem]{Remark}
\newcommand{\D}{\mathbb D}\newcommand{\R}{\mathbb R}
\newcommand{\C}{\mathbb C}\newcommand{\Qrat}{\mathbb Q}
\newcommand{\Adm}{\mathcal A}\newcommand{\F}{\mathcal F}
\newcommand{\T}{\mathcal T}\newcommand{\Kc}{\mathcal K}
\newcommand{\norm}[1]{\left\|#1\right\|_{\overline\D}}
\newcommand{\idzero}[1]{\mathbf1_{#1=0}}
\setlength{\parindent}{0pt}\setlength{\parskip}{4pt}
\setlength{\emergencystretch}{2em}
\pagestyle{myheadings}\markright{An exact algorithmic solution of the Belgian Chocolate Problem}
\title{\color{inkblue}An Exact Algorithmic Solution\newline of the Belgian Chocolate Problem}
\author{Xuandi Cheng}\date{Proof candidate v1.0.0 --- 13 September 2026}
\begin{document}\maketitle
\begin{abstract}
For $\delta>0$, the Belgian Chocolate Problem asks for nonzero real Hurwitz
stable polynomials $x,y,p$ with $\deg x\geq\deg y$ and
$p=(s^2-2\delta s+1)x+(s^2-1)y$. We specify an exact algorithm that
computes a critical real $Q$ by strict rational enclosures of any prescribed
precision and prove that the original admissible set is
$(0,(1-Q)/(1+Q))$. The analytic parameter set includes its endpoint $Q$,
but the strict polynomial parameter set does not. The proof combines
explicit uniform bounds for four analytic factors, a finite semialgebraic
hierarchy with an all-orders realization theorem, finite certificates on
both strict sides, and a two-test-point terminating subdivision algorithm.
Finite polynomial realization uses strict parameter slack; endpoint
exclusion uses an explicit polynomial perturbation that preserves actual
degrees even in the presence of leading-term cancellation. No degree bound,
closed-form value, or practical complexity estimate is assumed.
\end{abstract}

\section{Introduction}
This is a proof candidate submitted for independent verification, not a claim
of acceptance by the control-theory or complex-analysis communities.
The accompanying repository contains a kernel-checked concrete Lean main theorem;
its precise scope and its differences in proof implementation are stated in
Section~\ref{sec:formalcoverage}.
The problem considered here is the original polynomial simultaneous-
stabilization problem stated in Section~\ref{sec:original}. Its quantifiers
range over all finite degrees, and stability is strict. We determine its
entire admissible parameter set by specifying a real number through a
total certified approximation algorithm and proving the exact interval
criterion, including the excluded endpoint.

The proof distinguishes two issues. Open-disc analytic feasibility is
closed at its lower threshold; strictly stable polynomial feasibility is
open there. The analytic hierarchy therefore cannot be identified
pointwise with original polynomial feasibility at the threshold. The
precise interface is proved in Sections~\ref{sec:endpoint}--\ref{sec:return}.

The algorithm uses only integer and rational arithmetic, finite polynomial
enumeration, and effective real-algebraic decisions. Its output defines
the threshold independently of a feasibility query for the original
problem. Theorem~\ref{thm:main} is the final combined statement.

\section{Original simultaneous-stabilization problem}\label{sec:original}
\begin{definition}\label{def:original}
A nonzero $h\in\R[s]$ is \emph{Hurwitz stable} if every complex zero
$\lambda$ satisfies $\operatorname{Re}\lambda<0$. Nonzero constants are
stable. Put $a_\delta(s)=s^2-2\delta s+1$ and $b(s)=s^2-1$. Define
\begin{equation}\label{eq:original}
\Adm=\left\{\delta>0:\begin{array}{l}
 x,y,p\in\R[s]\setminus\{0\}\text{ are stable},\quad\deg x\geq\deg y,\\
 p=a_\delta x+by
\end{array}\right\}.
\end{equation}
The symbol $p$ denotes the third polynomial; $z$ will denote a disc variable.
\end{definition}
There is no leading-coefficient sign or monic normalization. The actual
degree of $p$ can be smaller than $\deg x+2$.

\begin{lemma}\label{lem:belowone}
If $\delta\in\Adm$, then $0<\delta<1$.
\end{lemma}
\begin{proof}
At $\delta=1$, evaluating the identity at $s=1$ gives $p(1)=0$.
For $\delta>1$, the roots $r_\pm=\delta\pm\sqrt{\delta^2-1}$ of
$a_\delta$ satisfy $0<r_-<1<r_+$. Evaluation gives
$p(r_\pm)/y(r_\pm)=r_\pm^2-1$, with opposite signs. But $p/y$ is a
continuous nonzero real function on $(0,\infty)$, so it has constant sign.
Both cases contradict stability.
\end{proof}

\begin{lemma}[Parameter change]\label{lem:parameters}
On $I=(0,1)$ the maps
\begin{equation}\label{eq:parameters}
 q(\delta)=\frac{1-\delta}{1+\delta},\qquad
 \delta(q)=\frac{1-q}{1+q}
\end{equation}
are mutually inverse strictly decreasing bijections.
\end{lemma}
\begin{proof}
Substitution gives the inverse identities. Each map takes $I$ to $I$
and has derivative $-2/(1+x)^2<0$.
\end{proof}
Write $P(q)$ for $\delta(q)\in\Adm$, for $q\in I$.


\begin{theorem}[Main result, stated in advance]\label{thm:front}
There is a critical real $Q\in(0,1/2)$ and a certified algorithm which, for every
$m\in\mathbb N$, terminates and returns $L,R\in\Qrat$ with
\[
 L<Q<R,\qquad R-L<2^{-m}.
\]
With $\Delta=(1-Q)/(1+Q)$, the original set is $\Adm=(0,\Delta)$ and
$\Delta\notin\Adm$. For the analytic predicate defined in the next section,
on $0<q<1$ one has $A(q)\Longleftrightarrow Q\le q$ and
$P(q)\Longleftrightarrow Q<q$. In particular $A(Q)$ holds but $P(Q)$ does not.
\end{theorem}
\begin{proof}
This is the principal assertion of Theorem~\ref{thm:main}, proved below.
The algorithm specifying $Q$ is given in Section~\ref{sec:algorithm}.
\end{proof}

\section{Analytic reformulation}\label{sec:analytic}
Let $\D=\{z\in\C:|z|<1\}$. For $0<t<1$, let $\F_t$ consist of
holomorphic $f:\D\to\C$ such that $f(\overline z)=\overline{f(z)}$,
the only zero is the simple zero $0$, and the only $1$-points are
the two simple points $it,-it$. Put
\[
 A(q):\ \F_{\sqrt q}\ne\varnothing,\qquad
 \T=\{q\in I:A(q)\},\qquad q_* =\inf\T.
\]
Nonemptiness and a positive lower bound for $\T$ are proved below.
The symbol $q_*$ is used in proofs, not in the algorithm's instructions.

\begin{lemma}[Polynomial-to-analytic map]\label{lem:forward}
For every $q\in I$, $P(q)$ implies $A(q)$.
\end{lemma}
\begin{proof}
For a polynomial witness put
\[
 w=\frac{s-1}{s+1},\qquad s=\frac{1+w}{1-w},\qquad
 f(w)=\frac{b(s)y(s)}{p(s)}.
\]
The second map is a biholomorphism from $\D$ onto the right half-plane.
Direct algebra gives
\begin{equation}\label{eq:cayley}
 (s+1)^2w=b(s),\qquad
 (s+1)^2(w^2+q)=(1+q)a_{\delta(q)}(s).
\end{equation}
Stability makes $f$ holomorphic, real symmetric, and gives
$1-f=a_{\delta(q)}x/p$. The only right-half-plane zero of $b$ is the
simple zero $1$, corresponding to $w=0$. The two distinct roots
$\delta\pm i\sqrt{1-\delta^2}$ of $a_\delta$ are in that half-plane
and correspond to $w=\pm i\sqrt q$. Neither $x,y$ nor $p$ has a zero
there. The change of variable has nonzero derivative, so the specified
fibres remain simple and there are no additional fibres. The boundary
point $w=1$ is not used; no assertion about $\deg p$ is needed.
\end{proof}

\begin{lemma}[Four-factor representation]\label{lem:factors}
For $q\in I$, $A(q)$ is equivalent to the existence of four
real-symmetric holomorphic functions $u,v,U,V$ on $\D$ satisfying
\begin{equation}\label{eq:factors}
 zu+(z^2+q)v=1,\qquad uU=1,\qquad vV=1.
\end{equation}
\end{lemma}
\begin{proof}
From $f$ take $u=f/z$, $v=(1-f)/(z^2+q)$, $U=1/u$, $V=1/v$.
The quotients have removable singularities at the specified points;
the exact simple fibres make them nonvanishing everywhere on $\D$.
Conversely let $f=zu$. The inverse identities make $u,v$ nonvanishing.
Thus $f$ has just the simple zero $0$, and
$1-f=(z^2+q)v$ has exactly the two specified simple zeros. Symmetry is
preserved in both directions, equivalently all Taylor coefficients are real.
\end{proof}

\section{Uniform bounds and effective coefficient boxes}\label{sec:bounds}
\subsection{Classical analytic inputs}
Poincar\'e densities are normalized to curvature $-1$, so
$\rho_\D(z)=2/(1-|z|^2)$. We use the following standard theorems in
the indicated forms.
\begin{description}
\item[Schwarz--Pick.]
If $h:G\to H$ is holomorphic between hyperbolic plane domains, then
$\rho_H(h(z))|h'(z)|\leq\rho_G(z)$.
\item[Schottky--Hempel density estimate.]
For $Y=\C\setminus\{0,1\}$,
\[
 \rho_Y(w)\geq\frac1{2|w|(C+|\log|w||)},\qquad
 C=\frac{\Gamma(1/4)^4}{4\pi^2}.
\]
This is the weakened form of the usual estimate in this normalization.
The Gamma integral gives
$\Gamma(1/4)<4+e^{-1}<5$. Since $\pi>3$, $C<625/36<18<64$, and hence
\begin{equation}\label{eq:density}
 \rho_Y(w)\geq\frac1{2|w|(64+|\log|w||)}.
\end{equation}
\item[Maximum modulus, argument principle, and Cauchy estimate.]
For a function holomorphic near a closed disc, its maximum modulus is
bounded by its boundary maximum. If it has no boundary zero, the image
boundary's winding number about zero equals the number of interior zeros,
with multiplicity. A boundary bound $M$ at radius $r$ gives Taylor
coefficient bound $Mr^{-j}$.
\item[Gronwall.]
If $a\geq0$ is absolutely continuous on $[0,\ell]$ and
$a'\leq ba$ almost everywhere, with $b$ integrable, then
$a(s)\leq a(0)\exp(\int_0^s b(\tau)\,d\tau)$.
\end{description}
The density estimate and Schwarz--Pick apply to arbitrary holomorphic
maps omitting $0,1$; no covering or local-univalence condition is imposed.

\subsection{An explicit all-functions bound}
Fix $\beta=3/4$. For rational $0\leq r<1$, define
\begin{equation}\label{eq:bounddata}
\begin{split}
 R&=\frac{1+\max\{r,\beta\}}2,\qquad
 d=\tfrac12\min\{R-\beta,1-R\},\\
 J&=\left\lceil16R/d\right\rceil+1,\qquad B=3^{66\cdot9^J},\\
 M(r)&=\left\lceil B\left(\frac1R+\frac1{R^2-\beta^2}
                         +R+R^2+\beta^2\right)\right\rceil.
\end{split}
\end{equation}
All denominators are positive and these are effective rational or integer
data, independent of $q_*$ and of all function degrees.

\begin{lemma}[Uniform four-factor bound]\label{lem:uniform}
For $0<t\leq\beta$, $f\in\F_t$, and its four factors,
\begin{equation}\label{eq:uniform}
 \max\{|u(z)|,|v(z)|,|U(z)|,|V(z)|\}\leq M(r)
 \qquad(|z|\leq r).
\end{equation}
\end{lemma}
\begin{proof}
On the circle $\gamma(\theta)=Re^{i\theta}$ the image $f\circ\gamma$
has winding numbers $1,2$ about $0,1$, respectively. It meets
\[
 K=\{w:1/4\leq|w|\leq2,\ |w-1|\geq1/4\}.
\]
Otherwise its connected image lies in one of the three disjoint regions
$|w|<1/4$, $|w-1|<1/4$, $|w|>2$. The first two force one winding
number to vanish. In the third the winding numbers about $0,1$ agree,
since the segment $[0,1]$ avoids the curve. Each is a contradiction.
Thus some $z_0$ on the circle satisfies
\begin{equation}\label{eq:anchor}
 1/4\leq|f(z_0)|\leq4,\qquad 1/4\leq|1-f(z_0)|\leq4.
\end{equation}

For every center $c$ on this circle, the disc of radius $d$ is contained
in $\D\setminus\{0,it,-it\}$: all deleted points and the unit boundary
are at distance at least $2d$. There $h=f$ and $h=1-f$ both omit $0,1$.
Schwarz--Pick and~\eqref{eq:density} give
\[
 \frac{|h'(z)|}{|h(z)|}
 \leq2(64+|\log|h(z)||)\frac{2d}{d^2-|z-c|^2}.
\]
Along a radial segment of length at most $d/2$, $|\log|h||$ is
absolutely continuous. The derivative inequality and Gronwall imply
\begin{equation}\label{eq:propagation}
 64+|\log|h(z)||\leq9(64+|\log|h(c)||),
 \qquad |z-c|\leq d/2,
\end{equation}
because
$\exp(2\int_0^{d/2}2d/(d^2-s^2)\,ds)=\exp(2\log3)=9$.
The circumference $2\pi R<8R$ can be traversed from $z_0$ in at most
$J$ steps of chord length at most $d/2$. Applying~\eqref{eq:propagation}
at their successive centers, using $\log4<2$ and $e<3$, yields
\[
 |f|,\ |1/f|,\ |1-f|,\ |1/(1-f)|\leq B\quad(|z|=R).
\]
Since $|z^2+t^2|\geq R^2-\beta^2$ there, the respective factor bounds
on this circle are
\begin{equation}\label{eq:fourbounds}
 B/R,\quad B/(R^2-\beta^2),\quad RB,\quad(R^2+\beta^2)B.
\end{equation}
All four factors are holomorphic on the entire disc. Apply the maximum
modulus principle separately and use their sum in~\eqref{eq:bounddata}.
\end{proof}

\subsection{Rational coefficient boxes}
For $m\geq1$ and $j\geq0$, put
\begin{equation}\label{eq:boxes}
 r_m=1-2^{-m},\qquad M_m=\max\{1,M(r_m)\},\qquad
 C_j=\min_{1\leq m\leq j+1}M_m r_m^{-j}.
\end{equation}
\begin{lemma}[Coefficient bounds and growth]\label{lem:coeff}
Each of the four Taylor coefficients at index $j$ has absolute value
at most $C_j$, for every $0<t\leq\beta$ and $f\in\F_t$. Moreover
\begin{equation}\label{eq:growth}
 \limsup_{j\to\infty}C_j^{1/j}\leq1.
\end{equation}
\end{lemma}
\begin{proof}
Cauchy's estimate gives each bound $M_m r_m^{-j}$ and hence their finite
minimum. For fixed $m$, $C_j\leq M_mr_m^{-j}$ when $j\geq m-1$.
The upper root growth is at most $r_m^{-1}$. Let $m\to\infty$.
\end{proof}
In particular every coefficient sequence in these boxes, not just one
already known to arise from a feasible function, defines a holomorphic
power series on $\D$.

\section{Finite semialgebraic hierarchy}\label{sec:hierarchy}
\begin{definition}\label{def:hierarchy}
For $q\in[0,9/16]$ and $N\geq0$, $H_N(q)$ is feasibility of the
following system in $4(N+1)$ real variables
$(u_j,v_j,U_j,V_j)_{0\leq j\leq N}$:
\begin{align}
 |u_j|,|v_j|,|U_j|,|V_j|&\leq C_j\quad(0\leq j\leq N),\label{eq:Hbox}\\
 u_{k-1}+v_{k-2}+qv_k&=\idzero{k}\quad(0\leq k\leq N),\label{eq:Hlinear}\\
 \sum_{i=0}^k u_iU_{k-i}=\sum_{i=0}^k v_iV_{k-i}
 &=\idzero{k}\quad(0\leq k\leq N).\label{eq:Hproduct}
\end{align}
Negative indices mean zero, and $\idzero{k}$ is one if $k=0$ and zero
otherwise. Absolute-value bounds mean the two corresponding linear
inequalities. No constraint is imposed on an infinite tail by this definition.
\end{definition}

\begin{proposition}[Finite exact decisions]\label{prop:decide}
For rational $q$ in the defined range, $H_N(q)$ is a finite rational
semialgebraic feasibility question, decidable by a terminating exact
procedure. Restricting a solution at level $N+1$ gives one at level $N$.
Also $H_0(0)$ is infeasible.
\end{proposition}
\begin{proof}
The data $C_j$ are rational, and all displayed constraints are polynomial
equalities or inequalities in finitely many variables. Effective
quantifier elimination for real closed fields decides such formulas.
Restriction retains precisely the lower-index equations. At $q=0$ the
zeroth linear equation is $0=1$.
\end{proof}
We use the standard effective real-algebraic theorem in its algorithmic
form: a fixed exact implementation, for example cylindrical algebraic
decomposition with real-root isolation, terminates and correctly decides
each finite rational first-order real-closed-field sentence. Merely a
nonconstructive choice of a truth value is not the algorithm used here.

\section{All-orders realization and the analytic threshold}\label{sec:realization}
\begin{theorem}[All-orders equivalence]\label{thm:allorders}
For $0<q\leq9/16$,
\begin{equation}\label{eq:allorders}
 A(q)\quad\Longleftrightarrow\quad\forall N\geq0\ H_N(q).
\end{equation}
\end{theorem}
\begin{proof}
Necessity follows from Lemmas~\ref{lem:factors} and~\ref{lem:coeff}.
For sufficiency take the compact product
\[
 \Kc=\prod_{j=0}^{\infty}[-C_j,C_j]^4.
\]
Countable products of compact intervals are compact in the coordinate
topology. Let $E_N(q)$ be the subset imposing only equations
\eqref{eq:Hlinear}--\eqref{eq:Hproduct} through index $N$. These are
closed cylinder sets, because each equation is continuous in finitely
many coordinates, and $E_{N+1}(q)\subseteq E_N(q)$. A finite feasible
solution, with all remaining coordinates set to zero, proves
$E_N(q)\ne\varnothing$. This extension is used only for nonemptiness
of that cylinder, not to satisfy higher equations.

The nested compact-set theorem gives a common infinite coefficient
object in $\bigcap_NE_N(q)$. By~\eqref{eq:growth} its four power series
converge absolutely on $\D$, locally uniformly, and are holomorphic
there. Cauchy products recover~\eqref{eq:factors} coefficient by
coefficient. In particular $uU=vV=1$ gives exact nonvanishing, rather
than a limiting non-strict inequality. Lemma~\ref{lem:factors} recovers
the simple fibres and real symmetry, and proves $A(q)$.
\end{proof}
No surjectivity of projections between feasible solution spaces is used.
The finite solutions are not chosen greedily to be compatible.

\begin{lemma}[Explicit strict seed]\label{lem:seed}
There are real rational polynomial factors at $q=1/2$ that are zero-free
on $\overline\D$. Consequently $\T$ is nonempty and $q_*<1/2$.
\end{lemma}
\begin{proof}
Take
\[
 u=-\frac{1000}{1089}(1+\tfrac45z)^2(1+\tfrac{17}{20}z),\qquad
 v=\frac2{1089}(1089+1000z+272z^2).
\]
Expanding $u=-(1000+2450z+2000z^2+544z^3)/1089$ verifies
$zu+(z^2+1/2)v=1$. The roots of $u$ are $-5/4,-20/17$.
The discriminant for $v$ is $-184832<0$, and its conjugate roots have
modulus $\sqrt{1089/272}>1$. Both polynomials remain nonvanishing on
some disc $|z|<R$ with $R>1$. For $f=zu$, $f(Rz)$ then belongs to
$\F_{\sqrt{1/2}/R}$, establishing a feasible parameter below $1/2$.
\end{proof}

\begin{lemma}[Scaling with strict slack]\label{lem:scale}
If $0<r<q<1$ and $A(r)$, then $A(q)$. In fact there are real-symmetric
holomorphic nonvanishing factors $u,v$ on $|z|<\sqrt{q/r}>1$ satisfying
$zu+(z^2+q)v=1$ there.
\end{lemma}
\begin{proof}
Take $f_0\in\F_{\sqrt r}$, $c=\sqrt{r/q}$, and $F(z)=f_0(cz)$.
On $|z|<1/c$ this has exactly the simple zero $0$ and the simple
$1$-points $\pm i\sqrt q$. The quotients $F/z$ and
$(1-F)/(z^2+q)$ are holomorphic nonvanishing factors on that disc.
\end{proof}

\begin{theorem}[Attained analytic threshold]\label{thm:analyticinterval}
\begin{equation}\label{eq:analyticinterval}
 0<q_*<1/2,\qquad \T=[q_*,1).
\end{equation}
\end{theorem}
\begin{proof}
The strict seed gives nonemptiness and the upper bound. At a feasible
$0<q\leq9/16$, the zeroth identity and box give
$qv_0=1$, $|v_0|\leq C_0$, and hence $q\geq1/C_0$.
A sequence of feasible parameters tending to the infimum can be chosen
eventually below $1/2$. Thus $q_*\geq1/C_0>0$.
Scaling proves upward closure, and the defining property of the infimum
then puts every $q>q_*$ in $\T$.

Choose feasible $q_n\to q_*$ eventually in $(0,1/2)$. Their four
coefficient sequences lie in the common compact metrizable space $\Kc$.
In $[0,1/2]\times\Kc$ take a convergent subsequence. Its parameter is
$q_*$; every finite coefficient equation passes to the limit by
continuity. The power-series construction in the proof of
Theorem~\ref{thm:allorders} gives four factors at $q_*$. Since $q_*>0$,
the two $1$-points stay distinct. Lemma~\ref{lem:factors} gives $A(q_*)$.
\end{proof}

\section{Finite certificates on both strict sides}\label{sec:certificates}
\begin{corollary}[Finite obstruction]\label{cor:negative}
For $0<q\leq9/16$, $q<q_*$ if and only if some $H_N(q)$ is infeasible.
For rational $q$ this is a finite, exact all-degrees obstruction.
\end{corollary}
\begin{proof}
Negate~\eqref{eq:allorders} and use~\eqref{eq:analyticinterval}.
For rational $q$, apply Proposition~\ref{prop:decide}. By
Lemma~\ref{lem:forward}, analytic impossibility also excludes any original
polynomial witness.
\end{proof}

\begin{definition}[Positive certificate]\label{def:positive}
For rational $q>0$, a positive certificate is $v\in\Qrat[z]$ with
$v(0)=1/q$ for which
\begin{equation}\label{eq:certificate}
 u(z)=\frac{1-(z^2+q)v(z)}z
\end{equation}
and $v$ are zero-free on $\overline\D$. The constant condition makes
$u$ a rational polynomial and the factor identity exact.
\end{definition}
This is decidable by real algebra: for $h\in\Qrat[z]$, write
$h(a+ib)=H_1(a,b)+iH_2(a,b)$. Closed-disc nonvanishing is
\[
 \neg\exists a,b\in\R:\ a^2+b^2\leq1,\ H_1(a,b)=H_2(a,b)=0.
\]
It contains only finite rational polynomial conditions.

\begin{lemma}[Exact polynomial approximation]\label{lem:approx}
If $0<r<q<1$ and $A(r)$, there are real polynomials $u_*,v_*$,
zero-free on $\overline\D$, with $zu_*+(z^2+q)v_*=1$.
When $q$ is rational they can be chosen rational.
\end{lemma}
\begin{proof}
Use the larger-disc factors from Lemma~\ref{lem:scale}. Their minimum
moduli $m_u,m_v$ on $\overline\D$ are positive. Choose
$0<\eta<\frac12\min\{m_v,m_u/(1+q)\}$. Taylor truncation approximates
$v$ within $\eta/2$ on $\overline\D$. For rational $q$, approximate
its finitely many nonconstant real coefficients by rationals with sum
of errors less than $\eta/2$, retaining $v_*(0)=1/q$ exactly. For real
$q$, retain the real Taylor polynomial. In either case
$\norm{v_*-v}<\eta$. Define $u_*$ by~\eqref{eq:certificate}, not by
independent approximation. Its numerator vanishes at zero. On $|z|=1$,
$|u_*-u|\leq(1+q)\norm{v_*-v}$. The difference is holomorphic at zero,
so the maximum modulus theorem gives this bound throughout the disc.
Both resulting polynomials therefore have positive minimum modulus.
The identity, symmetry, and rationality assertions are exact.
\end{proof}
The margins here depend on the individual strict pair $r<q$; no uniform
estimate as $q\downarrow q_*$ is used. The existence proof does not
require an algorithm to recover $f_0$ or its minimum moduli.

\begin{theorem}[Positive-certificate characterization]\label{thm:positive}
For rational $0<q\leq1/2$, a positive certificate exists if and only
if $q>q_*$.
\end{theorem}
\begin{proof}
Given a certificate, its two polynomial factors are nonvanishing on
some disc of radius $R>1$. Rescaling $f=zu$ to $f(Rz)$ gives
$A(q/R^2)$, and hence $q_*\leq q/R^2<q$.
Conversely for $q>q_*$ choose $r\in[q_*,q)$ with $A(r)$ and apply
Lemma~\ref{lem:approx}. The same existence argument works for every
real $q\in(q_*,1)$, requiring rational coefficients only when $q$ is
rational.
\end{proof}

\section{Certified computation of the critical real}\label{sec:algorithm}
Fix one deterministic exact real-closed-field decision implementation
as in Proposition~\ref{prop:decide}. For $s\geq0$, let $\mathcal E_s(q)$
be the finite list of rational polynomials of degree at most $s$, constant
term $1/q$, and other coefficients in
\[
 \{a/b:a\in\mathbb Z,\ |a|\leq s,\ 1\leq b\leq\max\{1,s\}\}.
\]
Order by actual degree, then increasing lexicographic order of rational
coefficient tuples, removing duplicates. Every rational polynomial
with that constant term occurs for some finite $s$.

Start at $[L_0,R_0]=[0,1/2]$, which strictly brackets $q_*$ by
Theorem~\ref{thm:analyticinterval}. Given a current bracket $[L,R]$, put
$a=(2L+R)/3$ and $b=(L+2R)/3$ and execute the following staged search.
\begin{quote}\ttfamily\small
for s = 0, 1, 2, ...:\\
\quad 1. Decide H\_s(a); if infeasible, return [a,R].\\
\quad 2. Decide H\_s(b); if infeasible, return [b,R].\\
\quad 3. Test E\_s(a) in its fixed order;\\
\qquad\ if a positive certificate is found, return [L,a].\\
\quad 4. Test E\_s(b) in its fixed order;\\
\qquad\ if a positive certificate is found, return [L,b].
\end{quote}
This implements parallel search by fair finite stages. Each stage
contains only finitely many terminating decisions.

\begin{theorem}[Total subdivision]\label{thm:step}
For rational $0\leq L<q_*<R\leq1/2$, the search terminates. Its output
is a nested strict bracket of $q_*$, with width at most $2(R-L)/3$.
\end{theorem}
\begin{proof}
The tests $a,b$ are distinct positive rationals below $1/2$, so at most
one equals $q_*$. The other has either a finite negative certificate
by Corollary~\ref{cor:negative}, or a finite positive certificate by
Theorem~\ref{thm:positive}. Exhaustiveness of the lists and the finite
staging ensure discovery at some finite stage. No earlier stage can
run forever. Negative certificates only raise the lower endpoint to a
point strictly below $q_*$. Positive certificates only lower the upper
endpoint to a point strictly above $q_*$. The four possible widths are
$2/3,1/3,1/3,2/3$ times the old width, respectively.
\end{proof}

\begin{theorem}[Algorithmic real and strict precision]\label{thm:Q}
Iteration gives a fully specified sequence of nested rational brackets
\begin{equation}\label{eq:brackets}
 L_j<q_*<R_j,\qquad R_j-L_j\leq\tfrac12(2/3)^j.
\end{equation}
Each stage is computable in finite time. Define $Q$ to be the unique
point of $\bigcap_j[L_j,R_j]$. Then $Q=q_*$. For each $m\geq0$ a
terminating algorithm returns rational $L_m^{\rm prec},R_m^{\rm prec}$
such that
\begin{equation}\label{eq:precision}
 L_m^{\rm prec}<Q<R_m^{\rm prec},\qquad
 R_m^{\rm prec}-L_m^{\rm prec}<2^{-m}.
\end{equation}
\end{theorem}
\begin{proof}
Induction using Theorem~\ref{thm:step} proves~\eqref{eq:brackets}.
Nested closed intervals with widths tending to zero have a unique common
point; since each contains $q_*$, it is $q_*$. For a precision input $m$,
choose the least $j$ with $\tfrac12(2/3)^j<2^{-m}$ by finite rational
comparisons, and return $[L_j,R_j]$. The geometric decay guarantees
termination of this index search, and every subdivision is terminating.
\end{proof}
The stage instructions use no $q_*$, $Q$, or original-feasibility oracle.
Their correctness proof identifies the independently defined $Q$ with
the analytic infimum. Fixing a different correct real-algebraic engine
does not change that limiting real. Define
\begin{equation}\label{eq:Delta}
 \Delta=\frac{1-Q}{1+Q}.
\end{equation}

\section{Endpoint separation}\label{sec:endpoint}
\begin{lemma}[Fixed actual degree stability]\label{lem:fixeddegree}
For each nonzero stable $h$ of actual degree $d$, there exists
$\tau_h>0$ such that every real polynomial $g$ of degree at most $d$
whose coefficients differ from those of $h$ by less than $\tau_h$
has degree exactly $d$ and is stable.
\end{lemma}
\begin{proof}
For $d=0$, use $|h(0)|/2$. Otherwise take disjoint closed circular discs
around the distinct roots, all in the open left half-plane. Let $\Gamma$
be their boundaries and
\[
 \mu=\min_\Gamma|h|>0,\quad B_h=\sum_{j=0}^d\max_\Gamma|s|^j,
 \quad\tau_h=\min\{|h_d|/2,\mu/(2B_h)\}.
\]
The leading coefficient stays nonzero, and $|g-h|<|h|$ on each circle.
Rouch\'e's theorem, for holomorphic functions near each closed disc,
preserves the number of enclosed zeros with multiplicity. These counts
total $d$; the fixed degree leaves no other roots. Repeated roots are
allowed.
\end{proof}

\begin{theorem}[Strict original-parameter perturbation]\label{thm:perturb}
If $\delta\in\Adm$, some $\delta'\in\Adm$ satisfies
$\delta<\delta'<1$, even when top coefficients of $a_\delta x+by$ cancel.
\end{theorem}
\begin{proof}
By Lemma~\ref{lem:belowone}, $0<\delta<1$. Fix a witness and put
$n=\deg x$, $m=\deg y\leq n$.
If $m<n$, then $\deg p=n+2$ with leading coefficient that of $x$.
Keep $x,y$ fixed and set $\delta'=\delta+e$,
$p'=p-2esx$. The identity is exact and the degree remains $n+2$.
If $M_x>0$ is the maximum absolute coefficient of $x$, any sufficiently
small $0<e<\min\{(1-\delta)/2,\tau_p/(4M_x)\}$ works by
Lemma~\ref{lem:fixeddegree}.

If $m=n$, let $e>0$ and define
\begin{equation}\label{eq:mixparams}
 k=\sqrt{1+e^2+2e/\delta},\qquad
 \delta'=\frac{k\delta}{1+\delta e},\qquad\phi(s)=ks-e.
\end{equation}
Use mixing coefficients
\begin{equation}\label{eq:mixcoeff}
\begin{aligned}
 \alpha_e&=1+e^2+e(\delta+1/\delta),&\beta_e&=e^2+e/\delta,\\
 \gamma_e&=e(1/\delta-\delta),&\theta_e&=1+e/\delta.
\end{aligned}
\end{equation}
Comparison of quadratic coefficients gives
\begin{equation}\label{eq:mixids}
 a_\delta(\phi)=\alpha_e a_{\delta'}+\gamma_e b,
 \qquad b(\phi)=\beta_e a_{\delta'}+\theta_e b.
\end{equation}
All coefficient checks are explicitly
\[
\begin{array}{lll}
 \alpha_e+\gamma_e=k^2,&\alpha_e-\gamma_e=1+e^2+2\delta e,
 &\alpha_e\delta'=k(\delta+e),\\
 \beta_e+\theta_e=k^2,&\beta_e-\theta_e=e^2-1,
 &\beta_e\delta'=ke.
\end{array}
\]
Define
\begin{equation}\label{eq:mixwitness}
 X=\alpha_e x(\phi)+\beta_e y(\phi),\quad
 Y=\gamma_e x(\phi)+\theta_e y(\phi),\quad P_1=p(\phi).
\end{equation}
Then $P_1=a_{\delta'}X+bY$ exactly. As $e\downarrow0$ the coefficients
of $X,Y,P_1$ tend to those of $x,y,p$. For all $e$ the first two degrees
are at most $n$; their leading coefficients tend to the original nonzero
ones. The actual degree of $P_1=p(ks-e)$ is exactly $\deg p$, since
$k>0$, regardless of how many leading terms have canceled originally.
Apply Lemma~\ref{lem:fixeddegree} at degree bounds $n,n,\deg p$.
For sufficiently small $e>0$ all three polynomials are stable, with
$\deg X=\deg Y=n$. One may take $e=2^{-j}$ for any sufficiently large
$j$ so that all three finite coefficient deviations are below their
positive tolerances. Thus the perturbation is genuinely nonzero.

Finally direct algebra gives
\begin{equation}\label{eq:direction}
 1-(\delta')^2=\frac{1-\delta^2}{(1+\delta e)^2}.
\end{equation}
For $e>0$ this lies strictly between $0$ and $1-\delta^2$, so
$\delta<\delta'<1$. Stability is asserted only for the small $e$
chosen above. Real coefficients, actual degrees, and the original
degree inequality are preserved without normalization.
\end{proof}

\begin{remark}[Why actual degree matters]
At $\delta=1/2$, $x=s+3/2$, $y=-s-1/2$, $p=s/2+2$ form a stable
witness with nominal degree three and actual degree one. Changing only
$\delta$ to $1/2+e$ gives
$-2es^2+(1/2-3e)s+2$, whose root product is $-1/e<0$ for $e>0$.
It has a positive root. The composition in~\eqref{eq:mixwitness}
preserves the actual degree instead; no root at infinity is disregarded.
\end{remark}

\begin{theorem}[Endpoint exclusion]\label{thm:endpoint}
For $q\in I$, $P(q)$ implies $q>Q$. In particular $\neg P(Q)$.
\end{theorem}
\begin{proof}
Perturb its original polynomial witness to $\delta'>\delta(q)$ by
Theorem~\ref{thm:perturb}. The corresponding $q'$ lies strictly between
$0$ and $q$. Lemma~\ref{lem:forward} gives $A(q')$, whence
$Q=q_*\leq q'<q$. This is an actual polynomial perturbation before
the infimum inequality is used, not an exclusion deduced from the
definition of an infimum alone.
\end{proof}

\section{Return to original polynomials}\label{sec:return}
\begin{lemma}[Cayley polynomial realization]\label{lem:lift}
Let $q\in I$ and let real polynomials $u,v$ be zero-free on
$\overline\D$, with $wu+(w^2+q)v=1$. Put $n=\max\{\deg u,\deg v\}$.
Then
\begin{equation}\label{eq:lift}
 \begin{split}
 x(s)&=(1+q)(s+1)^n v\left(\frac{s-1}{s+1}\right),\\
 y(s)&=(s+1)^n u\left(\frac{s-1}{s+1}\right),\qquad
 p(s)=(s+1)^{n+2}
 \end{split}
\end{equation}
are nonzero real stable polynomials with $\deg x=\deg y=n$ and
$p=a_{\delta(q)}x+by$. Rational input data give rational coefficients.
\end{lemma}
\begin{proof}
For $h(w)=h_d\prod_{j=1}^d(w-\omega_j)$ with $d\leq n$,
\begin{equation}\label{eq:rootlift}
 (s+1)^n h\left(\frac{s-1}{s+1}\right)
 =h_d(s+1)^{n-d}\prod_{j=1}^d((1-\omega_j)s-(1+\omega_j)).
\end{equation}
Zero-freeness gives $|\omega_j|>1$, so $\omega_j\ne1$ and
\[
 \operatorname{Re}\frac{1+\omega_j}{1-\omega_j}
 =\frac{1-|\omega_j|^2}{|1-\omega_j|^2}<0.
\]
Any additional roots are $-1$. Empty products include constants.
The coefficient of $s^n$ is $h(1)\ne0$, proving the actual degree $n$.
Apply this to $u,v$ and use~\eqref{eq:cayley} to obtain
\[
 a_{\delta(q)}x+by=(s+1)^{n+2}((w^2+q)v+wu)=(s+1)^{n+2}.
\]
The equality, first computed for $s\ne-1$, holds as a polynomial
identity. Denominator clearing uses only real, or rational, operations.
\end{proof}

\begin{theorem}[Off-endpoint realization]\label{thm:offendpoint}
If $0<r<q<1$ and $A(r)$, then $P(q)$. Consequently every real
$q\in(Q,1)$ is polynomial feasible, with a rational polynomial witness
whenever $q$ is rational.
\end{theorem}
\begin{proof}
Apply Lemma~\ref{lem:approx} and then Lemma~\ref{lem:lift}.
For $q>Q$ take $r=Q$, using analytic attainment. No uniform estimate
in $q-Q$ is needed. This applies on all of $I$, not just the range of
the finite hierarchy.
\end{proof}

\section{Main theorem}\label{sec:main}
\begin{theorem}[Exact algorithmic solution]\label{thm:main}
The algorithm of Section~\ref{sec:algorithm} uniquely determines a real
$Q\in(0,1/2)$ and terminates for every requested strict precision
\eqref{eq:precision}. With $\Delta=(1-Q)/(1+Q)$, the original problem
has exactly the admissible set
\begin{equation}\label{eq:main}
 \boxed{\Adm=(0,\Delta),\qquad \Delta\notin\Adm,\qquad1/3<\Delta<1.}
\end{equation}
For every $q\in(0,1)$ the two predicates are precisely
\begin{equation}\label{eq:separation}
 \boxed{A(q)\Longleftrightarrow q\geq Q,\qquad
        P(q)\Longleftrightarrow q>Q.}
\end{equation}
Thus for $0<q\leq9/16$ the all-layer analytic equivalence remains valid
including $Q$, but its polynomial counterpart is the strict-slack relation
\begin{equation}\label{eq:slackhierarchy}
 P(q)\Longleftrightarrow\exists r\in(0,q)\ A(r)
 \Longleftrightarrow\exists r\in(0,q)\ \forall N\ H_N(r).
\end{equation}
The first equivalence in~\eqref{eq:slackhierarchy} holds on all of $I$.
\end{theorem}
\begin{proof}
Theorem~\ref{thm:Q} proves the algorithmic assertions and identifies
$Q=q_*$. Theorem~\ref{thm:analyticinterval} gives the first predicate
in~\eqref{eq:separation}. Theorems~\ref{thm:endpoint} and
\ref{thm:offendpoint} give the second. Lemma~\ref{lem:belowone} excludes
$\delta\geq1$; the decreasing inverse change of parameter gives
$q(\delta)>Q$ exactly when $0<\delta<\Delta$. The bounds on $\Delta$
follow from $0<Q<1/2$. Strict slack follows either from the perturbation
and forward map, or the explicit intervals; realization proves its
reverse direction. For $q\leq9/16$, every $r\in(0,q)$ is in the
range of Theorem~\ref{thm:allorders}.
\end{proof}

\begin{corollary}[Relative closure and endpoint meaning]\label{cor:closure}
In $I=(0,1)$, the analytic feasible set $[Q,1)$ is the relative closure
of the polynomial feasible set $(Q,1)$. At $Q$ all $H_N(Q)$ are feasible,
but no original polynomial witness exists. No endpoint factors can extend
nonvanishingly to a disc of radius $R>1$.
\end{corollary}
\begin{proof}
The set and hierarchy assertions follow from the preceding theorems.
If endpoint factors extended as stated, their identity would extend by
the identity theorem, and $Rz\,u(Rz)$ would belong to
$\F_{\sqrt Q/R}$, contradicting $Q=\inf\T$. Open-disc analytic
attainment therefore does not supply a closed-disc polynomial certificate.
The closure in $\R$ instead includes $1$, corresponding to $\delta=0$
outside the original domain. In $(0,\infty)$ the original set has
relative closure $(0,\Delta]$.
\end{proof}

\section{Exactness, computability, and non-stabilization}\label{sec:discussion}
\subsection{No finite level stabilizes at the threshold}
Put $S_N=\{q\in[0,1/2]:H_N(q)\}$ and $\alpha_N=\min S_N$.
\begin{theorem}[Strict finite-level non-stabilization]\label{thm:nonstable}
Each $S_N$ is nonempty compact semialgebraic and $\alpha_N$ is an
effectively computable real algebraic number. Moreover
\begin{equation}\label{eq:nonstable}
 0<\alpha_N<Q,\qquad \alpha_N\uparrow Q,\qquad
 \forall N\ \exists M>N:\ \alpha_M>\alpha_N.
\end{equation}
For every fixed $N$ there is $\varepsilon_N>0$ such that $H_N(q)$
holds whenever $|q-Q|<\varepsilon_N$. The neighborhood can be chosen
inside the defined parameter range.
\end{theorem}
\begin{proof}
The joint constraints in $q$ and the finite coefficient box are closed
and bounded, so their projection is compact. The seed at $1/2$ makes
$S_N$ nonempty. Quantifier elimination and real-root isolation give its
algebraic minimum effectively. The zeroth equation gives
$\alpha_N\geq1/C_0>0$.
Restriction gives $S_{N+1}\subseteq S_N$. All-orders realization and the
analytic interval give $\bigcap_NS_N=[Q,1/2]$. If $\alpha$ is the
increasing limit of the minima, then for each fixed $N$ its tail lies
in the closed set $S_N$, so $\alpha\in S_N$. Hence $\alpha=Q$.

For strictness, take endpoint factors $u_*,v_*,U_*,V_*$. The four
positive individual bounds in~\eqref{eq:fourbounds} each lie strictly
below their sum used in $M(r)$. Cauchy's estimate and the finite minimum
in~\eqref{eq:boxes} therefore give strict inequalities for each true
coefficient: $|u_{*,j}|,|v_{*,j}|,|U_{*,j}|,|V_{*,j}|<C_j$.
For $q$ near $Q$ define analytic germs at zero by
\[
 v_q=v_*+1/q-1/Q,\qquad u_q=\frac{1-(z^2+q)v_q}{z}.
\]
The quotient is removable, $v_q(0)=1/q$, and
$u_q(0)=(q/Q)u_*(0)\ne0$. Their inverses are holomorphic germs, whose
every fixed finite collection of coefficients depends continuously on
$q$. The three identities hold for the germs. At $Q$ the finitely many
box constraints have strict margin, so continuity keeps them inside
their boxes on some neighborhood. Thus $H_N$ has feasible points below
$Q$, proving $\alpha_N<Q$. Convergence then forces a later strict
increase for every fixed $N$.
\end{proof}
The germs in this proof need not have zero-free inverses on the whole
disc. No single $\varepsilon$ or infinite box control is asserted.
One can also search effectively for a later $H_M(\alpha_N)$ that is
infeasible, using exact real-algebraic input representation and the
all-orders obstruction. Non-stabilization does not affect the precision
algorithm, and asserts nothing about other possible formulas for $Q$.

\subsection{Exact determination and limits of membership algorithms}
The unique interval-intersection definition of $Q$ uses the complete
specified program and its proved convergence, not a finite decimal
approximation. Together with both directions of~\eqref{eq:main} and
the independent endpoint exclusion, it determines the original feasible
set exactly. Absence of a closed form, an algebraicity result, or a useful
complexity bound does not weaken these statements.
The computed brackets transform to
\[
 \frac{1-R_j}{1+R_j}<\Delta<\frac{1-L_j}{1+L_j},\qquad
 \frac{1-L_j}{1+L_j}-\frac{1-R_j}{1+R_j}
 =\frac{2(R_j-L_j)}{(1+L_j)(1+R_j)}\leq2(R_j-L_j).
\]
One extra bit of strict precision for $Q$ suffices for any requested
strict precision for $\Delta$.

For an effectively given real input promised different from $\Delta$,
successively refining the input and threshold intervals eventually
decides the strict comparison. For rational $\delta>0$, $\delta\geq1$
is impossible and $\delta\leq1/3$ is feasible. In the remaining range
the rational transform lies in $(0,1/2)$ and the two certificate searches
decide all nonendpoint inputs. The parameter explicitly specified as
$\Delta$ is infeasible by the endpoint theorem. No irrationality of
$Q$ or $\Delta$ is claimed, and no total rational-input membership
algorithm at a potentially rational endpoint is asserted.

This must be distinguished from the unconditional totality of the
precision algorithm: every precision input terminates, with no
nonendpoint promise. A general algorithm deciding membership from an
arbitrary black-box real Cauchy name would face the discontinuity of the
interval characteristic function at its endpoint; it is not a conclusion
of this paper. The exact interval theorem and endpoint exclusion do not
require such an equality detector.

\section{Scope and limitations}\label{sec:scope}
Here \emph{exact algorithmic solution} means a complete feasible-set theorem
together with arbitrary prescribed precision, strict rational enclosures,
and proved correctness and finite termination of the precision procedure.
It does not mean a closed form. No algebraicity, minimal polynomial,
special-function closed form, practical complexity bound, fixed finite
truncation deciding $Q$, or generic finite-time black-box equality test
$q=Q$ is claimed. No efficient numerical implementation or benchmarked
approximation of $Q$ is supplied in this release. A precision algorithm
and a total endpoint-equality decision procedure are different assertions.

\section{Formal verification and reproducibility}\label{sec:formalcoverage}
The \href{https://github.com/chengxuandi/belgian-chocolate-problem-lean-proof-20260910210030-init}{public repository}
contains the source, verification scripts and tagged release.
The release identifier is \texttt{v1.0.0-proof-candidate}. It uses Lean 4.33.0
and Mathlib 4.33.0 (the dependency commit is pinned in
\texttt{lake-manifest.json}). The concrete main declaration is
\begin{center}\small
\path{BelgianChocolate.belgianChocolate_exact_algorithmic_solution}
\end{center}
in \texttt{BCPThreshold/FinalMain.lean}. Its statement has no pending
analytic, certificate, endpoint, or realization-interface hypothesis.
The exact parameter sets, endpoint separation, all-orders equivalence
on $0<q\le9/16$, and arbitrary-precision enclosure assertion are included.

This verifies the main conclusion and its actual Lean dependency graph,
not every paper assertion verbatim. The coefficient proof uses the proved
coarse Schottky/cosine-lift development rather than the displayed Hempel
density theorem, reaching the same boxes $C_j$. The negative search uses
finite compact-polynomial infeasibility certificates with proved completeness,
not formalized general CAD. The positive verifier also uses finite polynomial
data. These implement the same two-sided mechanism but need not reproduce
the stage-by-stage outputs of Section~\ref{sec:algorithm}; their strict
brackets identify the same real $Q$.

The fixed-level algebraic-minimum and non-stabilization discussion of
Section~\ref{sec:discussion} is a paper-level result; it is not a separately
exported theorem of this Lean release, and is not needed by the final Lean
main theorem. The total mathematical iterator uses \texttt{Nat.find} and
is marked \texttt{noncomputable}; its finite-prefix searches are executable,
and their finite success and the iteration bounds are proved. No separate
native threshold-computation application is claimed.

The main theorem and concrete interface bundle have only
\texttt{propext}, \texttt{Classical.choice}, and \texttt{Quot.sound}
as transitive axiom dependencies. The BCP source has no project axiom
declarations or proof holes. Modeling fidelity still requires human review.

With \texttt{elan}, Git, Python 3 and Bash installed, clone the repository,
run \texttt{lake exe cache get}, then \texttt{./verify\_bcp.sh}.
The script builds the library and every BCP module, runs the final transitive
axiom audit, and checks source for proof holes and axiom declarations.
The verification guide gives the exact commands and the short reviewer
route. Build warnings about style or deprecated APIs are not proof failures.

\section{AI assistance and independent review}
AI models assisted mathematical exploration, route search, adversarial
review, proof engineering, and Lean formalization. The manuscript and Lean
source are made public for independent examination. Kernel checking does
not replace human review of the modeling choices, mathematical exposition,
novelty, or scope of the claim. Independent verification by control-theory,
complex-analysis, and formal-methods experts is actively requested.

\end{document}
